\documentclass[11pt, letterpaper]{article}

% ── Paquetes ──────────────────────────────────────────────────────────────────
\usepackage[utf8]{inputenc}
\usepackage[T1]{fontenc}
\usepackage[spanish]{babel}
\usepackage{geometry}
\usepackage{enumitem}
\usepackage{xcolor}
\usepackage{hyperref}
\usepackage{fancyhdr}
\usepackage{titlesec}
\usepackage{booktabs}
\usepackage{graphicx}
\usepackage{parskip}

% ── Configuración de página ───────────────────────────────────────────────────
\geometry{margin=2.5cm}

% ── Colores institucionales IMSS ──────────────────────────────────────────────
\definecolor{imssgreen}{HTML}{006847}
\definecolor{imssgray}{HTML}{4A4A4A}
\definecolor{alertorange}{HTML}{D4760A}
\definecolor{lightgreen}{HTML}{E8F5E9}
\definecolor{lightgray}{HTML}{F5F5F5}

% ── Configuración de hyperref ─────────────────────────────────────────────────
\hypersetup{
    colorlinks=true,
    linkcolor=imssgreen,
    urlcolor=imssgreen
}

% ── Estilos de sección ────────────────────────────────────────────────────────
\titleformat{\section}
    {\large\bfseries\color{imssgreen}}{\thesection.}{0.5em}{}
    [\vspace{-0.5em}\textcolor{imssgreen}{\rule{\textwidth}{0.4pt}}]

% ── Encabezado y pie de página ────────────────────────────────────────────────
\pagestyle{fancy}
\fancyhf{}
\fancyhead[L]{\small\color{imssgray}EpiForecast-MX}
\fancyhead[R]{\small\color{imssgray}Avance 3 --- Retroalimentación}
\fancyfoot[C]{\small\color{imssgray}\thepage}
\renewcommand{\headrulewidth}{0.4pt}

% ── Comando para caja de comentario ──────────────────────────────────────────
\newcommand{\comentario}[2]{%
    \vspace{0.5em}
    \noindent
    \fcolorbox{alertorange!40}{lightgray}{%
        \begin{minipage}[t]{0.95\textwidth}
            \textbf{\color{alertorange}Comentario #1.}~#2
        \end{minipage}
    }
    \vspace{0.5em}
}

% ── Comando para reconocimiento positivo ──────────────────────────────────────
\newcommand{\felicitacion}[1]{%
    \vspace{0.5em}
    \noindent
    \fcolorbox{imssgreen!40}{lightgreen}{%
        \begin{minipage}[t]{0.95\textwidth}
            \textbf{\color{imssgreen}Reconocimiento.}~#1
        \end{minipage}
    }
    \vspace{0.5em}
}

% ══════════════════════════════════════════════════════════════════════════════
\begin{document}

% ── Título ────────────────────────────────────────────────────────────────────
\begin{center}
    {\Large\bfseries\color{imssgreen} Retroalimentación del Avance~3 --- Baseline}\\[0.4em]
    {\large EpiForecast-MX: Pronóstico Epidemiológico para el IMSS}\\[1.2em]
    \begin{tabular}{rl}
        \textbf{Asesora:}          & Dra.\ Grettel Barceló Alonso \\
        \textbf{Fecha:}            & 19 de febrero de 2026 \\
        \textbf{Intento:}          & 1 \\
    \end{tabular}
\end{center}

\vspace{1em}

% ── Reconocimiento general ────────────────────────────────────────────────────
\section{Reconocimiento General}

\felicitacion{%
Es excelente ver que, además del análisis de los modelos, hayan implementado un
dashboard de métricas, lo que aporta un valor significativo a la presentación y
comprensión de los resultados. Como en cada entrega, ustedes van más allá de lo
solicitado. ¡Felicitaciones!}

% ── Comentarios detallados ────────────────────────────────────────────────────
\section{Comentarios Detallados}

\comentario{1}{%
\textbf{Horizonte de validación cruzada vs.\ horizonte de pronóstico.}
¿Por qué el horizonte de la validación cruzada no coincide con el horizonte del
pronóstico (52 semanas)? Con esto las métricas pueden dar una estimación
demasiado optimista.}

\comentario{2}{%
\textbf{Proximidad entre métricas y umbrales aceptables.}
Presentar las métricas del modelo (Sección~3) separadas de los umbrales
aceptables (Sección~7) dificulta que el lector pueda evaluar rápidamente si los
resultados son buenos o malos. Recomiendo mover ``Desempeño mínimo aceptable''
justo después de \emph{Resultados}.}

\comentario{3}{%
\textbf{Especificación de la columna de brecha.}
Cuando incluyan la brecha deben especificar con qué métrica se está calculando.
Usar \texttt{Brecha (MAPE pp)}. Mi recomendación es que omitan esta columna
porque ya la incluyen en la Sección~5 \emph{Análisis de Sub/Sobreajuste}. La
tabla de métricas muestra datos objetivos, mientras que la brecha es
interpretación.}

\comentario{4}{%
\textbf{Pie de figura incorrecto (Figura~10).}
No es correcto el pie de figura. Sería bueno incluir a qué estado se refieren
los atípicos.}

\comentario{5}{%
\textbf{Reportar parámetros del modelo Prophet.}
Al usar Prophet, es importante reportar los parámetros del modelo
(\texttt{changepoint\_prior\_scale}, \texttt{seasonality\_prior\_scale},
\texttt{seasonality\_mode}, \ldots), incluso si son los predeterminados.}

\comentario{6}{%
\textbf{Horizonte de predicción vs.\ estacionalidad anual.}
Dado que el horizonte de predicción utilizado (24 semanas) es más corto que la
periodicidad de la estacionalidad principal (anual), la evaluación del modelo no
refleja completamente la capacidad de capturar la estacionalidad anual.}

\comentario{7}{%
\textbf{Acentos en títulos de gráficos.}
Asegúrense de poner acentos en los títulos de los gráficos, no sólo en el
reporte, sino también en el dashboard.}

% ── Resumen de acciones ──────────────────────────────────────────────────────
\section{Resumen de Acciones Requeridas}

\begin{enumerate}[label=\textbf{A\arabic*.}, leftmargin=2em, itemsep=0.4em]
    \item Alinear el horizonte de validación cruzada con el horizonte de
          pronóstico de 52 semanas.
    \item Reorganizar el reporte: mover la sección de desempeño mínimo aceptable
          inmediatamente después de los resultados.
    \item Eliminar la columna de brecha de la tabla de métricas (o, en su
          defecto, especificar \texttt{Brecha (MAPE pp)}).
    \item Corregir el pie de la Figura~10 e indicar el estado al que
          corresponden los valores atípicos.
    \item Incluir una tabla con los hiperparámetros de Prophet utilizados.
    \item Evaluar la extensión del horizonte de predicción a 52 semanas para
          capturar la estacionalidad anual completa.
    \item Revisar y corregir los acentos en todos los títulos de gráficos del
          reporte y del dashboard.
\end{enumerate}

\vspace{2em}
\noindent\textcolor{imssgray}{\rule{\textwidth}{0.4pt}}
\begin{center}
    \small\color{imssgray}
    Documento generado para seguimiento interno del equipo EpiForecast-MX.
\end{center}

\end{document}
